\pdfvariable suppressoptionalinfo 611
\providecommand{\CRevisionRound}{2}
\documentclass[9pt]{extarticle}
\usepackage[a4paper,margin=0.62in]{geometry}
\usepackage{fontspec}
\setmainfont{Latin Modern Roman}
\setsansfont{Latin Modern Sans}
\setmonofont{Latin Modern Mono}
\usepackage[english]{babel}
\usepackage{amsmath,amssymb,booktabs,microtype}
\usepackage[hidelinks]{hyperref}
\emergencystretch=1.5em
\setlength{\parindent}{1em}
\setlength{\parskip}{0.20ex}
\newcommand{\R}{\mathbb R}
\newcommand{\sech}{\operatorname{sech}}
\title{The Focusing Cubic NLS Soliton:\\Exact Hessian Spectrum and Pöschl--Teller Factorization}
\author{Route-A structural certificate HCS-C221}
\date{}

\begin{document}
\maketitle
\begin{abstract}
We close a convention-locked one-dimensional theorem for the focusing cubic
nonlinear Schrödinger equation.  The sech standing wave has exact mass,
Hamiltonian, action and Vakhitov--Kolokolov (VK) slope.  Its two real Hessians have
the complete Pöschl--Teller discrete spectrum, essential threshold, kernels and
Morse index.  The zero-frequency and defocusing faces are separated, and no
finite-box spectrum or arithmetic interpretation is used.
\ifnum\CRevisionRound>0
The variational scaling and explicit eigenfunctions are included.
\fi
\ifnum\CRevisionRound>1
The ladder factorization, boundary ledger, independent evidence counts and
strict Route-A stop are included.
\fi
\end{abstract}
\noindent\textbf{Keywords:} nonlinear Schrödinger; soliton; Hessian;
Pöschl--Teller; Morse index; Route A.

\section{Equation and standing wave}
We study
\begin{equation}
 i\psi_t+\psi_{xx}+2|\psi|^2\psi=0,
 \qquad \psi(t,x)=e^{i\omega t}Q_\omega(x),\qquad \omega>0,
\end{equation}
on \(\R\).  Direct differentiation of
\(Q_\omega(x)=\sqrt\omega\sech(\sqrt\omega x)\) gives
\begin{equation}
 -Q_\omega''+\omega Q_\omega-2Q_\omega^3=0.
\end{equation}
For
\[
 M(u)=\int_\R|u|^2dx,\quad
 H(u)=\tfrac12\int_\R|u_x|^2dx-\tfrac12\int_\R|u|^4dx,
 \quad S_\omega=H+\tfrac\omega2M,
\]
the substitution \(y=\sqrt\omega x\), \(z=\tanh y\) yields
\begin{equation}
 M(Q_\omega)=2\sqrt\omega,\quad
 \|Q_\omega'\|_2^2=\tfrac23\omega^{3/2},\quad
 \|Q_\omega\|_4^4=\tfrac43\omega^{3/2},
\end{equation}
and hence
\begin{equation}
 H(Q_\omega)=-\tfrac13\omega^{3/2},\qquad
 S_\omega(Q_\omega)=\tfrac23\omega^{3/2},\qquad
 \frac{d}{d\omega}M(Q_\omega)=\omega^{-1/2}>0.
\end{equation}

\section{The complete Hessian spectrum}
The real and imaginary second variations are the self-adjoint operators
\begin{equation}
 L_+=-\partial_x^2+\omega-6\omega\sech^2(\sqrt\omega x),\qquad
 L_-=-\partial_x^2+\omega-2\omega\sech^2(\sqrt\omega x).
\end{equation}
The theorem is
\begin{equation}
 \sigma_{\rm ess}(L_+)=\sigma_{\rm ess}(L_-)=[\omega,\infty),
 \quad \sigma_{\rm disc}(L_+)=\{-3\omega,0\},
 \quad \sigma_{\rm disc}(L_-)=\{0\}.
\end{equation}
The \(-3\omega\) eigenfunction is \(\sech^2(\sqrt\omega x)\); the zero
eigenfunctions are \(Q_\omega'\) for \(L_+\) and \(Q_\omega\) for \(L_-\).
All listed eigenvalues are simple.  Thus \(L_+\) has Morse index one and
kernel \(\operatorname{span}\{Q_\omega'\}\), while \(L_-\ge0\) and has kernel
\(\operatorname{span}\{Q_\omega\}\).

To see the signs without a box discretization, put \(y=\sqrt\omega x\) and
\[
 A_\ell=\partial_y+\ell\tanh y,\qquad
 A_\ell^*=-\partial_y+\ell\tanh y.
\]
Multiplication gives
\begin{equation}
 -\partial_y^2+1-6\sech^2y=A_2^*A_2-3,
 \qquad -\partial_y^2+1-2\sech^2y=A_1^*A_1.
\end{equation}
The \(\ell=1,2\) ladder has exactly the bound states displayed in (6); the
free remainder starts at one, which rescales to the threshold \(\omega\).

\ifnum\CRevisionRound>0
\section{Boundaries and evidence}
At \(\omega\downarrow0\), the profile and discrete eigenvalues collapse to the
essential threshold.  Reversing the cubic sign has no nonzero bright \(H^1\)
sech branch.  A periodic domain gives elliptic/cnoidal waves, and dimensions
\(d\ge2\) have different criticality; neither is silently included.
\fi
\ifnum\CRevisionRound>1
The executable receipt contains 15 profile, 3 integral, 15 spectral and 15
factorization rows.  An independent checker passes 497 assertions, SymPy
passes 19 identities, replay is byte exact, and 17 repaired/stale-hash hostile
mutations are rejected.  These rows are regression evidence only; they do not
prove the continuum statement by sampling.  In particular, no full nonlinear
orbital-stability theorem is asserted.

\section{Strict Route-A boundary}
The continuum NLS coefficients carry no intrinsic rational-prime or
prime-power clock, and the standing-wave symmetry supplies only a weak relative
orbit hint.  The strict tuple is
\[
 (A0,A1,A2,A3,A4)=\texttt{(A0\_FAIL,A1\_WEAK,A2\_FAIL,A3\_FAIL,}
 \texttt{A4\_NATURAL\_QUANTIZATION)}.
\]
The result is \texttt{ROUTE\_A\_REJECTED}; Route B is disabled.  The final
entry records only a natural Hamiltonian quantization hint, not a
Hilbert--Pólya operator.  Hessian traces and continuous spectra are not
dynamical zeta functions, Euler factors or target divisors.
The locked scope literal is
\texttt{NO\_BAD\_EULER\_OR\_ROOT\_NUMBER}.
The source baseline lock for this round is
\texttt{86c7bb8a39cdd1b8e941e45833b068170ca06287}.

The source-attributed references are Zakharov--Shabat
\href{https://www.jetp.ras.ru/cgi-bin/dn/e_034_01_0062}{[ZS72]}, Weinstein
\href{https://doi.org/10.1137/0516034}{[Weinstein85]}, and the historical
Pöschl--Teller solvable-potential/ladder source
\href{https://doi.org/10.1007/BF01331132}{[PT33]}; no priority claim is made.
\fi

\ifcase\CRevisionRound
\paragraph{Revision focus.} Round 0 freezes (1)--(2) and the operator owner.
\or
\paragraph{Revision focus.} Round 1 adds (3)--(6), explicit kernels and the
zero-frequency/defocusing boundary.
\or
\paragraph{Revision focus.} Round 2 adds factorization (7), all scope faces,
independent evidence counts, provenance and the strict Route-A stop.
\fi

\begin{thebibliography}{3}
\footnotesize
\setlength{\itemsep}{0pt}
\bibitem{ZS72} V. E. Zakharov and A. B. Shabat, ``Exact theory of two-dimensional
self-focusing and one-dimensional self-modulation of waves in nonlinear media,''
\emph{Soviet Physics JETP} 34 (1972), 62--69.
\href{https://www.jetp.ras.ru/cgi-bin/dn/e_034_01_0062}{Official record}.
\bibitem{Weinstein85} M. I. Weinstein, ``Modulational Stability of Ground States
of Nonlinear Schrödinger Equations,'' \emph{SIAM J. Math. Anal.} 16 (1985),
472--491. DOI: \href{https://doi.org/10.1137/0516034}{10.1137/0516034}.
\bibitem{PoschlTeller1933} G. Pöschl and E. Teller, ``Bemerkungen zur
Quantenmechanik des anharmonischen Oszillators,'' \emph{Zeitschrift für Physik}
83 (1933), 143--151. DOI:
\href{https://doi.org/10.1007/BF01331132}{10.1007/BF01331132}.  Cited only as
the historical solvable Pöschl--Teller potential/ladder source.
\end{thebibliography}

\section*{Declarations}
\textbf{Data and code.} Synthetic exact probes and package-local deterministic
code accompany the manuscript; no training data are used.\quad
\textbf{Ethics/funding.} No human or animal data; no external funding reported.\quad
\textbf{AI disclosure.} An AI coding assistant supported drafting and exact-code
development; it was not an external reviewer or independent error process.

\end{document}
